\chapter{One-pagers and drills}
\label{app:drills}

\section{The cards}

Five documents in \code{onepagers/}, each building to one A4 side, meant to be
printed and on the desk rather than read on a screen.

\begin{kittable}{@{}llX@{}}
\textbf{Card} & \textbf{Chapter} & \textbf{When} \\
\midrule
\code{run-sheet-card} & \ref{ch:runsheet} & beside the keyboard during a build round \\
\code{decision-method-card} & \ref{ch:answering} & any technical conversation \\
\code{failure-tree-card} & \ref{ch:breaks} & beside the run sheet \\
\code{pre-call-briefing-card} & \ref{ch:person} & filled in per call, read beforehand \\
\code{banned-claims-card} & \ref{ch:claims} & filled in per process, read beforehand \\
\end{kittable}

The first three are printed once. The last two are filled in each time, which
is why they are a card rather than a page of this book.

\judgement{Print them. A card on the desk is consulted under pressure and a
file is not, and the two that matter most during a build round are the run
sheet and the failure tree \dash{} both of which exist to be looked at exactly
when you are least able to recall their contents.}

\section{The rehearsal}

The protocol is Chapter~\ref{ch:rehearsal}. The materials are in
\code{drills/}.

\code{drills/README.md} carries the protocol and the instructions for whoever
is playing the owner. \code{drills/scoring.md} carries the table with targets.
\code{drills/role-briefs/} carries four briefs in four unrelated domains.

\begin{kittable}{@{}lX@{}}
\textbf{Brief} & \textbf{The domain} \\
\midrule
\code{01-bakery-orders} & a bakery taking cake orders by telephone \\
\code{02-clinic-scheduling} & a physiotherapy clinic booking appointments \\
\code{03-warehouse-picking} & a small warehouse picking and packing \\
\code{04-venue-ticketing} & a music venue selling tickets \\
\end{kittable}

Choose one you know nothing about. The point of the exercise is that you have
to ask, and a domain you understand lets you skip the part being measured.

\section{How the briefs are built}

Each has four parts, and the third is the instrument:

\textbf{What the owner knows} and will say when asked directly.

\textbf{How they behave} \dash{} vague about volumes, specific about
irritations, uninterested in technology.

\textbf{What they do not say until asked outright.} Every brief contains a
second actor and a constraint that is never volunteered. A candidate who does
not ask builds the wrong thing, which is exactly what a real session would have
done to them.

\textbf{Four timed interjections} they hold and you do not know about: two
requirement changes and two challenges, at fixed elapsed times, one of them
after the build has started.

\judgement{The mid-build interjection is the most informative minute of the
rehearsal. What you do when a requirement changes at minute sixty-five is the
thing a real session will find out about you, and most people have never
watched themselves do it.}
